\chapter{绪论}
根据需要，作者也可将本部分命名为“引言”。本章应说明论文的主题和选题范围，评述相关研究，并介绍研究内容及章节安排。本文保留原模板的结构与写作提示，示例内容不代表实际研究成果。
\section{研究背景及意义}
说明研究问题、应用背景、理论意义和实践价值。正文采用中文宋体，英文与数字采用Times New Roman，小四号，首行缩进两个汉字，1.5倍行距。
\section{国内外研究现状分析}
结合历史回顾、已有文献和理论分析，对研究现状进行综合评述。
\subsection{三级标题}
三级标题使用阿拉伯数字分级编号，编号末尾不加点号。需要引用文献时，使用引用命令自动生成标注。
\subsection{研究问题}
明确已有研究的不足及本论文拟解决的问题，避免将工作计划写成已经验证的结论。
\section{研究内容}
根据实际研究填写研究内容、技术路线和创新点。
\section{组织结构}
从第二章起依次介绍各章的主要内容。
